\documentclass[11pt,a4paper]{article}
\usepackage{fontspec,amsmath,amssymb,amsthm,microtype}
\setmainfont{lmroman10-regular.otf}[BoldFont=lmroman10-bold.otf,ItalicFont=lmroman10-italic.otf,BoldItalicFont=lmroman10-bolditalic.otf]
\usepackage{xcolor}
\usepackage[most]{tcolorbox}
\usepackage{mathrsfs}
\providecommand{\lt}{<}
\providecommand{\gt}{>}
\setmonofont{DejaVuSansMono.ttf}[BoldFont=DejaVuSansMono-Bold.ttf,ItalicFont=DejaVuSansMono-Oblique.ttf,BoldItalicFont=DejaVuSansMono-BoldOblique.ttf,Scale=0.86]
\newfontfamily\leanfallback{STIXTwoMath-Regular.otf}[Scale=0.86]
\definecolor{leanrule}{gray}{0.62}
\definecolor{leanback}{gray}{0.962}
\definecolor{leantitle}{gray}{0.22}
\definecolor{leanmuted}{gray}{0.38}
\newtcolorbox{leanbox}[1]{enhanced,breakable,colback=leanback,colframe=leanrule,
  boxrule=0.4pt,arc=0.6mm,left=2.4mm,right=2.4mm,top=1.8mm,bottom=1.8mm,
  before skip=0.8em,after skip=1em,fonttitle=\small\bfseries,coltitle=leantitle,
  colbacktitle=leanback,titlerule=0.3pt,title={#1}}
\newlength{\leancharwidth}
\newenvironment{leancode}{\par\smallskip\begingroup\ttfamily\footnotesize\raggedright
  \settowidth{\leancharwidth}{M}\setlength{\parskip}{0pt}\setlength{\parindent}{0pt}
  \frenchspacing}{\par\endgroup\smallskip}
\newcommand{\leanline}[2]{\hangindent=\dimexpr#1\leancharwidth+4\leancharwidth\relax
  \hangafter=1\noindent\hspace*{#1\leancharwidth}#2\par}
\newcommand{\leansource}[1]{\par{\footnotesize\color{leanmuted}#1}\par}
\newcommand{\leanlabel}[1]{\par\smallskip\noindent{\small\bfseries #1}\par}
\usepackage[margin=25mm]{geometry}
\usepackage[unicode,hidelinks]{hyperref}
\hypersetup{pdfauthor={LangevinLean Project}}
\setlength{\emergencystretch}{3em}
\setlength{\parskip}{0.4em}
\setlength{\parindent}{0pt}
\allowdisplaybreaks[2]
\raggedbottom
\title{\texorpdfstring{Finite smoothing comparison for mean-field integral equations}{Finite smoothing comparison for mean-field integral equations}}
\author{LangevinLean Project}
\date{}
\begin{document}\maketitle
Draft mathematical exposition; these notes may contain mistakes.


Assume additionally that the convex penalty is differentiable with Lipschitz gradient. Its proximal residual then equals the original gradient evaluated at the proximal point. The finite proximal displacement bounds the difference between original and envelope gradients by a constant times the smoothing scale. Construct both mean-field integral solutions with identical initial state and additive input. The input cancels on subtraction, and Grönwall bounds the supremum over time of their mean-square norms.




Earlier results used here: \cite[Section~3.43]{LangevinLean2026ULASC}.

\section{\texorpdfstring{Equality of the subgradient and gradient at a differentiability point}{Equality of the subgradient and gradient at a differentiability point}}


If \(s\) is a supporting subgradient, then \(x\mapsto f(x)-\langle s,x\rangle\) has a global minimum at \(p\). Differentiability makes its derivative zero, giving \(s=\nabla f(p)\). Applied to the proximal residual and envelope-gradient formula, this gives \(\nabla f^\eta(x)=\nabla f(P_\eta(x))\).



Equivalently, a support \(s\) at \(p\) minimizes \(f(x)-\langle s,x\rangle\) at \(p\), so \(s=\nabla f(p)\). The conclusion is an equality at the minimizing point, not an asymptotic approximation.

\begin{leanbox}{Lean: \texttt{IsSubgradient.\allowbreak{}eq\_\allowbreak{}gradient}, \texttt{gradient\_\allowbreak{}moreau\_\allowbreak{}eq\_\allowbreak{}gradient\_\allowbreak{}prox}}
\leansource{\texttt{IsSubgradient.\allowbreak{}eq\_\allowbreak{}gradient} \quad ProximalLangevin/\allowbreak{}MeanFieldSmoothing.lean, lines 29–43}
\begin{leancode}
\leanline{0}{theorem\ IsSubgradient.eq\_gradient\ \{f\ :\ E\ →\ ℝ\}\ \{p\ s\ :\ E\}}
\leanline{4}{(hs\ :\ IsSubgradient\ f\ p\ s)\ (hf\ :\ DifferentiableAt\ ℝ\ f\ p)\ :\ s\ =\ gradient\ f\ p}
\end{leancode}
\leansource{\texttt{gradient\_\allowbreak{}moreau\_\allowbreak{}eq\_\allowbreak{}gradient\_\allowbreak{}prox} \quad ProximalLangevin/\allowbreak{}MeanFieldSmoothing.lean, lines 46–52}
\begin{leancode}
\leanline{0}{theorem\ gradient\_moreau\_eq\_gradient\_prox\ \{f\ :\ E\ →\ ℝ\}}
\leanline{4}{(hconv\ :\ ConvexOn\ ℝ\ univ\ f)\ (hb\ :\ ∃\ b,\ ∀\ z,\ b\ ≤\ f\ z)}
\leanline{4}{(hf\ :\ Differentiable\ ℝ\ f)\ \{η\ :\ ℝ\}\ (hη\ :\ 0\ <\ η)\ (x\ :\ E)\ :}
\leanline{4}{gradient\ (moreau\ f\ η)\ x\ =\ gradient\ f\ (prox\ f\ η\ x)}
\end{leancode}
\leanlabel{How the Lean proof follows the argument}
Formal correspondence: IsSubgradient.eq\_gradient : A supporting subgradient of a differentiable function is its actual gradient. gradient\_moreau\_eq\_gradient\_prox : The actual Moreau score is the original smooth score at the true proximal point.



\texttt{Is\allowbreak{}Subgradient.\allowbreak{}eq\_\allowbreak{}gradient}: the intermediate facts include \texttt{F}, \texttt{hmin}, \texttt{h}, \texttt{hd}, \texttt{hg}, in source order; the source steps apply or rewrite \texttt{inner\_\allowbreak{}sub\_\allowbreak{}right}, \texttt{map\_\allowbreak{}sub}, \texttt{Pi.\allowbreak{}sub\_\allowbreak{}def}, \texttt{Inner\allowbreak{}Product\allowbreak{}Space.\allowbreak{}to\allowbreak{}Dual\_\allowbreak{}apply\_\allowbreak{}apply}, \texttt{sub\_\allowbreak{}eq\_\allowbreak{}zero.\allowbreak{}mp}. The remaining scalar rearrangements are checked by \texttt{linarith}.



\texttt{gradient\_\allowbreak{}moreau\_\allowbreak{}eq\_\allowbreak{}gradient\_\allowbreak{}prox}: the source steps apply or rewrite \texttt{gradient\_\allowbreak{}moreau}, \texttt{subgradient\_\allowbreak{}of\_\allowbreak{}is\allowbreak{}Prox\allowbreak{}Point}, \texttt{prox\_\allowbreak{}is\allowbreak{}Prox\allowbreak{}Point}, \texttt{eq\_\allowbreak{}gradient}.
\end{leanbox}

\section{\texorpdfstring{The finite gradient difference under smoothing}{The finite gradient difference under smoothing}}


The gradient \(\nabla f\) is \(L_2\)-Lipschitz, and \(P_\eta\) is nonexpansive. Their composition therefore has constant \(L_2\) independently of \(\eta\). The displacement bound \(\|x-P_\eta(x)\|\le\eta K\) then gives

\[\|\nabla f(x)-\nabla f^\eta(x)\|\le L_2\|x-P_\eta(x)\|\le\eta L_2K.\]

Optimality and differentiability identify \(\nabla f^\eta(x)=\nabla f(P_\eta x)\). When \(\nabla f\) is \(L_2\)-Lipschitz and \(f\) is \(K\)-Lipschitz, the same argument gives

\[\|\nabla f^\eta(x)-\nabla f(x)\|\le L_2\|P_\eta x-x\|\le\eta L_2K.\]

The envelope gradient is \(\nabla f\circ P_\eta\). Its gradient constant is bounded by \(L_2\), which can be smaller than the generic \(1/\eta\) bound. Comparison with \(\nabla f(x)\) costs at most \(L_2\|P_\eta x-x\|\le\eta L_2K\). This finite uniform defect supplies the forcing in the integral-equation comparison.

\begin{leanbox}{Lean: \texttt{smooth\_\allowbreak{}moreau\_\allowbreak{}gradient\_\allowbreak{}lipschitz}, \texttt{gradient\_\allowbreak{}moreau\_\allowbreak{}defect\_\allowbreak{}le}}
\leansource{\texttt{smooth\_\allowbreak{}moreau\_\allowbreak{}gradient\_\allowbreak{}lipschitz} \quad ProximalLangevin/\allowbreak{}MeanFieldSmoothing.lean, lines 55–63}
\begin{leancode}
\leanline{0}{theorem\ smooth\_moreau\_gradient\_lipschitz\ \{f\ :\ E\ →\ ℝ\}}
\leanline{4}{(hconv\ :\ ConvexOn\ ℝ\ univ\ f)\ (hb\ :\ ∃\ b,\ ∀\ z,\ b\ ≤\ f\ z)}
\leanline{4}{(hf\ :\ Differentiable\ ℝ\ f)\ \{η\ :\ ℝ\}\ (hη\ :\ 0\ <\ η)}
\leanline{4}{\{L\ :\ ℝ≥0\}\ (hLip\ :\ LipschitzWith\ L\ (gradient\ f))\ :}
\leanline{4}{LipschitzWith\ L\ (gradient\ (moreau\ f\ η))}
\end{leancode}
\leansource{\texttt{gradient\_\allowbreak{}moreau\_\allowbreak{}defect\_\allowbreak{}le} \quad ProximalLangevin/\allowbreak{}MeanFieldSmoothing.lean, lines 67–77}
\begin{leancode}
\leanline{0}{theorem\ gradient\_moreau\_defect\_le\ \{f\ :\ E\ →\ ℝ\}}
\leanline{4}{(hconv\ :\ ConvexOn\ ℝ\ univ\ f)\ (hb\ :\ ∃\ b,\ ∀\ z,\ b\ ≤\ f\ z)}
\leanline{4}{(hf\ :\ Differentiable\ ℝ\ f)\ \{η\ :\ ℝ\}\ (hη\ :\ 0\ <\ η)}
\leanline{4}{\{L\ K\ :\ ℝ≥0\}\ (hLip\ :\ LipschitzWith\ L\ (gradient\ f))\ (hK\ :\ LipschitzWith\ K\ f)\ (x\ :\ E)\ :}
\leanline{4}{‖gradient\ f\ x\ -\ gradient\ (moreau\ f\ η)\ x‖\ ≤\ η\ *\ (L\ :\ ℝ)\ *\ (K\ :\ ℝ)}
\end{leancode}
\leanlabel{How the Lean proof follows the argument}
Formal correspondence: smooth\_moreau\_gradient\_lipschitz : Smoothness of the original score is preserved by the actual Moreau envelope. gradient\_moreau\_defect\_le : The uniform score defect is derived from proximal optimality and displacement. Its coefficient is independent of the smoothing parameter.



\texttt{smooth\_\allowbreak{}moreau\_\allowbreak{}gradient\_\allowbreak{}lipschitz}: the intermediate fact is \texttt{heq}, in source order; the source steps apply or rewrite \texttt{mul\_\allowbreak{}one}, \texttt{h\allowbreak{}Lip.\allowbreak{}comp}, \texttt{prox\_\allowbreak{}lipschitz}.



\texttt{gradient\_\allowbreak{}moreau\_\allowbreak{}defect\_\allowbreak{}le}: the source steps apply or rewrite \texttt{gradient\_\allowbreak{}moreau\_\allowbreak{}eq\_\allowbreak{}gradient\_\allowbreak{}prox}. The \texttt{calc} block checks the successive equalities or inequalities in the argument. The remaining scalar rearrangements are checked by \texttt{ring}.
\end{leanbox}

\section{\texorpdfstring{Component regularity from the sum and the smooth term}{Component regularity from the sum and the smooth term}}


If \(U=f+g\) and \(f\) are differentiable, then \(g=U-f\) is differentiable. Gradient linearity gives \(\nabla g=\nabla U-\nabla f\). Taking differences at two points and the triangle inequality yields \(\operatorname{Lip}(\nabla g)\le L_U+L_f\).



Thus the constant \(L_2\) may be supplied from regularity of the original sum and the known smooth component. This gives a justified \(L_2\) without adding an unstated differentiability assumption on the remaining component.

\begin{leanbox}{Lean: \texttt{component\_\allowbreak{}differentiable\_\allowbreak{}of\_\allowbreak{}sum}, \texttt{component\_\allowbreak{}gradient\_\allowbreak{}lipschitz\_\allowbreak{}of\_\allowbreak{}sum}}
\leansource{\texttt{component\_\allowbreak{}differentiable\_\allowbreak{}of\_\allowbreak{}sum} \quad ProximalLangevin/\allowbreak{}MeanFieldSmoothing.lean, lines 82–88}
\begin{leancode}
\leanline{0}{theorem\ component\_differentiable\_of\_sum\ \{f\ g\ :\ E\ →\ ℝ\}}
\leanline{4}{(hfg\ :\ Differentiable\ ℝ\ (fun\ x\ =>\ f\ x\ +\ g\ x))\ (hf\ :\ Differentiable\ ℝ\ f)\ :}
\leanline{4}{Differentiable\ ℝ\ g}
\end{leancode}
\leansource{\texttt{component\_\allowbreak{}gradient\_\allowbreak{}lipschitz\_\allowbreak{}of\_\allowbreak{}sum} \quad ProximalLangevin/\allowbreak{}MeanFieldSmoothing.lean, lines 92–103}
\begin{leancode}
\leanline{0}{theorem\ component\_gradient\_lipschitz\_of\_sum\ \{f\ g\ :\ E\ →\ ℝ\}}
\leanline{4}{(hfg\ :\ Differentiable\ ℝ\ (fun\ x\ =>\ f\ x\ +\ g\ x))\ (hf\ :\ Differentiable\ ℝ\ f)}
\leanline{4}{\{LU\ Lf\ :\ ℝ≥0\}\ (hLU\ :\ LipschitzWith\ LU\ (gradient\ (fun\ x\ =>\ f\ x\ +\ g\ x)))}
\leanline{4}{(hLf\ :\ LipschitzWith\ Lf\ (gradient\ f))\ :}
\leanline{4}{LipschitzWith\ (LU\ +\ Lf)\ (gradient\ g)}
\end{leancode}
\leanlabel{How the Lean proof follows the argument}
Formal correspondence: component\_differentiable\_of\_sum : Original-model differentiability and smoothness derive the missing component regularity, rather than requiring a separate inverse-scale estimate. component\_gradient\_lipschitz\_of\_sum : If the full original score and first component are Lipschitz, the second component's score has their summed Lipschitz constant.



\texttt{component\_\allowbreak{}differentiable\_\allowbreak{}of\_\allowbreak{}sum}: the proof is the application \texttt{convert hfg.\allowbreak{}sub hf using 1}. It specializes that result with the hypotheses and objects in the exact statement.



\texttt{component\_\allowbreak{}gradient\_\allowbreak{}lipschitz\_\allowbreak{}of\_\allowbreak{}sum}: the intermediate facts include \texttt{hg}, \texttt{heq}, in source order; the source steps apply or rewrite \texttt{gradient\_\allowbreak{}add\_\allowbreak{}of\_\allowbreak{}differentiable}, \texttt{h\allowbreak{}LU.\allowbreak{}sub}.
\end{leanbox}

\section{\texorpdfstring{Comparison with common additive input}{Comparison with common additive input}}


Let integral solutions with drifts \(g,k\) share additive input \(B\) and have initial states \(x,y\). Assume one drift is globally Lipschitz and \(\sup_z\|g(z)-k(z)\|\le\varepsilon\). Subtraction cancels the input. Adding and subtracting one drift at the other state bounds the integrand by \(L\|Y-Z\|+\varepsilon\). Integral Grönwall gives

\[\|Y_t-Z_t\|\le e^{Lt}\|x-y\|+\varepsilon\Psi(L,t),\qquad\Psi(L,t)=\begin{cases}(e^{Lt}-1)/L,&L\ne0,\\t,&L=0.\end{cases}\]

No differentiation of \(B\) is needed. The bounded defect is integrated with factor \(\Psi\). The definition of \(\Psi\) includes zero Lipschitz constant, so the formula requires no division by zero.

\begin{leanbox}{Lean: \texttt{integral\_\allowbreak{}drift\_\allowbreak{}perturbation}}
\leansource{\texttt{integral\_\allowbreak{}drift\_\allowbreak{}perturbation} \quad ProximalLangevin/\allowbreak{}MeanFieldSmoothing.lean, lines 113–162}
\begin{leancode}
\leanline{0}{theorem\ integral\_drift\_perturbation}
\leanline{4}{(g\ k\ :\ H\ →\ H)\ (L\ M\ :\ ℝ≥0)\ (hg\ :\ LipschitzWith\ L\ g)\ (hk\ :\ LipschitzWith\ M\ k)}
\leanline{4}{(ε\ :\ ℝ)\ (hdefect\ :\ ∀\ x,\ ‖g\ x\ -\ k\ x‖\ ≤\ ε)}
\leanline{4}{(x\ y\ :\ H)\ (B\ Y\ Z\ :\ ℝ\ →\ H)\ (hY\ :\ Continuous\ Y)\ (hZ\ :\ Continuous\ Z)}
\leanline{4}{(T\ :\ ℝ)\ (hT\ :\ 0\ ≤\ T)}
\leanline{4}{(hYe\ :\ ∀\ t\ ∈\ Icc\ 0\ T,\ Y\ t\ -\ x\ =\ -(∫\ s\ in\ 0..t,\ g\ (Y\ s))\ +\ B\ t)}
\leanline{4}{(hZe\ :\ ∀\ t\ ∈\ Icc\ 0\ T,\ Z\ t\ -\ y\ =\ -(∫\ s\ in\ 0..t,\ k\ (Z\ s))\ +\ B\ t)\ :}
\leanline{4}{∀\ t\ ∈\ Icc\ 0\ T,\ ‖Y\ t\ -\ Z\ t‖\ ≤\ gronwallBound\ ‖x\ -\ y‖\ (L\ :\ ℝ)\ ε\ t}
\end{leancode}
\leanlabel{How the Lean proof follows the argument}
Formal correspondence: integral\_drift\_perturbation : Two actual integral equations driven by common additive input differ by the Grönwall response to a uniform drift defect. No derivative of the input is used, and no error recurrence is assumed.



\texttt{integral\_\allowbreak{}drift\_\allowbreak{}perturbation}: the intermediate facts include \texttt{D}, \texttt{F}, \texttt{h\allowbreak{}D}, \texttt{hbound}, \texttt{h}, in source order; the source steps apply or rewrite \texttt{sub\_\allowbreak{}eq\_\allowbreak{}iff\_\allowbreak{}eq\_\allowbreak{}add}, \texttt{has\allowbreak{}Deriv\allowbreak{}At\_\allowbreak{}const}, \texttt{dist\_\allowbreak{}eq\_\allowbreak{}norm\_\allowbreak{}sub}, \texttt{dist\_\allowbreak{}triangle}, \texttt{h\allowbreak{}Y.\allowbreak{}sub}. The \texttt{calc} block checks the successive equalities or inequalities in the argument.
\end{leanbox}

\section{\texorpdfstring{The finite smoothing estimate for the two integral systems}{The finite smoothing estimate for the two integral systems}}


Assume \(g_2\) is differentiable with \(L_2\)-Lipschitz gradient, and retain convexity, a lower bound, and value Lipschitz constant \(K\). The supporting residual equals the gradient at the proximal minimizer:

\[
\nabla g_2^\eta(v)=\nabla g_2(P_\eta(v)).
\]

Nonexpansiveness of \(P_\eta\) preserves constant \(L_2\), and \(\|v-P_\eta(v)\|\le\eta K\) gives

\[
\|\nabla g_2(v)-\nabla g_2^\eta(v)\|\le\eta L_2K.
\]

If instead regularity is stated for \(U=g_1+g_2\) with constant \(L_U\) and for \(g_1\) with constant \(L_1\), the difference argument supplies \(L_2=L_U+L_1\).



Construct both Hilbert-space solutions with common initial state and continuous additive input. The averaged drifts differ uniformly by at most \(2\eta L_2K\). Their pair-gradient constant is at most \(L_1+L_2\), independent of \(\eta\). The averaging operator gives the corresponding factor two. Grönwall yields

\[
\sup_{0\le t\le T}\|Y_t-Z_t\|_{L^2}
 \le 2\eta L_2K\,\Psi(2(L_1+L_2),T),
\qquad
\Psi(a,T)=\begin{cases}(e^{aT}-1)/a,&a>0,\\T,&a=0.\end{cases}
\]

For deterministic initial parameter and latent-only Brownian input these are the mean-field systems described above. The finite bound is linear in \(\eta\). More explicitly the averaged drift constant is \(2(L_1+L_2)\) and its uniform difference is \(2\eta L_2K\). The supremum is outside the \(L^2\) norm. Proposition C.3 instead bounds the expectation of the supremum of squared errors. The square root of that quantity bounds the supremum of the \(L^2\) error norms. A bound on the expectation of the supremum of squared errors is not established here.

\begin{leanbox}{Lean: \texttt{meanFieldSmoothingConstant}, \texttt{exists\_\allowbreak{}meanField\_\allowbreak{}smoothing\_\allowbreak{}comparison}}
\leansource{\texttt{meanFieldSmoothingConstant} \quad ProximalLangevin/\allowbreak{}MeanFieldSmoothing.lean, lines 169–172}
\begin{leancode}
\leanline{0}{def\ meanFieldSmoothingConstant\ (L₁\ L₂\ K\ :\ ℝ≥0)\ (T\ :\ ℝ)\ :\ ℝ\ :=}
\leanline{2}{2\ *\ (L₂\ :\ ℝ)\ *\ (K\ :\ ℝ)\ *\ gronwallBound\ 0\ (2\ *\ ((L₁\ :\ ℝ)\ +\ (L₂\ :\ ℝ)))\ 1\ T}
\leanline{0}{\strut}
\leanline{0}{section\ MeanField}
\end{leancode}
\leansource{\texttt{exists\_\allowbreak{}meanField\_\allowbreak{}smoothing\_\allowbreak{}comparison} \quad ProximalLangevin/\allowbreak{}MeanFieldSmoothing.lean, lines 187–249}
\begin{leancode}
\leanline{0}{theorem\ exists\_meanField\_smoothing\_comparison}
\leanline{4}{(initial\ :\ Lp\ (IPLA.Joint\ dθ\ dx)\ 2\ Q)}
\leanline{4}{(B\ :\ ℝ\ →\ Lp\ (IPLA.Joint\ dθ\ dx)\ 2\ Q)\ (hB\ :\ Continuous\ B)\ (hB0\ :\ B\ 0\ =\ 0)}
\leanline{4}{(T\ :\ ℝ)\ (hT\ :\ 0\ ≤\ T)\ :}
\leanline{4}{∃\ (hU\ :\ LipschitzWith\ (L₁\ +\ L₂)\ (gradient\ (fun\ x\ =>\ g₁\ x\ +\ g₂\ x)))}
\leanline{6}{(hV\ :\ LipschitzWith\ (L₁\ +\ L₂)\ (gradient\ (smoothedPotential\ g₁\ g₂\ η)))}
\leanline{6}{(Y\ Z\ :\ ℝ\ →\ Lp\ (IPLA.Joint\ dθ\ dx)\ 2\ Q),}
\leanline{6}{Continuous\ Y\ ∧\ Continuous\ Z\ ∧\ Y\ 0\ =\ initial\ ∧\ Z\ 0\ =\ initial\ ∧}
\leanline{6}{(∀\ t\ ∈\ Icc\ 0\ T,\ Y\ t\ -\ initial\ =}
\leanline{8}{-(∫\ s\ in\ 0..t,\ lpMeanFieldDrift\ Q\ (fun\ x\ =>\ g₁\ x\ +\ g₂\ x)\ hU\ (Y\ s))\ +\ B\ t)\ ∧}
\leanline{6}{(∀\ t\ ∈\ Icc\ 0\ T,\ Z\ t\ -\ initial\ =}
\leanline{8}{-(∫\ s\ in\ 0..t,\ lpMeanFieldDrift\ Q\ (smoothedPotential\ g₁\ g₂\ η)\ hV\ (Z\ s))\ +\ B\ t)\ ∧}
\leanline{6}{∀\ t\ ∈\ Icc\ 0\ T,\ ‖Y\ t\ -\ Z\ t‖\ ≤\ η\ *\ meanFieldSmoothingConstant\ L₁\ L₂\ K\ T}
\end{leancode}
\leanlabel{How the Lean proof follows the argument}
Both systems are constructed with the same continuous additive input. Subtracting their equations cancels that input and leaves a differentiable integral error. One drift contributes its global Lipschitz constant times the current difference; the other contributes the uniform Moreau score defect just proved. Grönwall then bounds the actual process difference by a coefficient times \(\eta\) on the finite horizon. The component-regularity declarations justify the smoothness constants of each summand from the stated assumptions on the regular sum before this perturbation argument is applied.



Formal correspondence: meanFieldSmoothingConstant : Explicit finite-horizon coefficient for the mean-field \(L^2\) smoothing error. The Grönwall response is \(T\) at zero smoothness and otherwise \(\frac{\exp(2(L_1+L_2)T)-1}{2(L_1+L_2)}\). exists\_meanField\_smoothing\_comparison : Constructed original and smoothed mean-field processes share their actual initial variable and continuous additive input. Their \(L^2\) distance, uniformly over the finite horizon, is bounded by \(\eta\) times an explicit \(\eta\)-independent constant. This is \(\sup_t\|\operatorname{error}_t\|_{L^2}\), not a pathwise supremum inside an expectation.



\texttt{mean\allowbreak{}Field\allowbreak{}Smoothing\allowbreak{}Constant} is a definition with defining expression \texttt{2 * (L\ensuremath{{}_{2}} : \ensuremath{\mathbb R}) * (K : \ensuremath{\mathbb R}) * gronwall\allowbreak{}Bound 0 (2 * ((L\ensuremath{{}_{1}} : \ensuremath{\mathbb R}) + (L\ensuremath{{}_{2}} : \ensuremath{\mathbb R}))) 1 T}. The right-hand side specifies the object used in the mathematical statement.



\texttt{exists\_\allowbreak{}mean\allowbreak{}Field\_\allowbreak{}smoothing\_\allowbreak{}comparison}: the intermediate facts include \texttt{h\allowbreak{}U}, \texttt{h\allowbreak{}V}, \texttt{U}, \texttt{heq}, \texttt{hpt}, in source order; the source steps apply or rewrite \texttt{gradient\_\allowbreak{}add\_\allowbreak{}of\_\allowbreak{}differentiable}, \texttt{gradient\_\allowbreak{}moreau\_\allowbreak{}defect\_\allowbreak{}le}, \texttt{norm\_\allowbreak{}zero}, \texttt{hpt.\allowbreak{}trans}, \texttt{htime.\allowbreak{}trans\_\allowbreak{}eq}. The remaining scalar rearrangements are checked by \texttt{ring}, \texttt{nlinarith}.
\end{leanbox}


Earlier results used here: \cite[Section~3.13.2]{LangevinLean2026ULASC}; \cite[Section~3.13.3]{LangevinLean2026ULASC}; \cite[Section~3.1.1]{LangevinLean2026IPLA}.

\begin{thebibliography}{ULA-LSI}
\bibitem[ULA-SC]{LangevinLean2026ULASC}
LangevinLean Project. \emph{Unadjusted Langevin under Strong Convexity}. LL-01, 2026. Draft report, content revision d34ead06987c. \url{https://odakyildiz.github.io/LangevinLean/docs/pdf/UnadjustedLangevin/FullProof.pdf}.
\bibitem[IPLA-SC]{LangevinLean2026IPLA}
LangevinLean Project. \emph{Interacting Particle Langevin under Strong Convexity}. LL-03, 2026. Draft report, content revision ec0201df2207. \url{https://odakyildiz.github.io/LangevinLean/docs/pdf/IPLA/FullProof.pdf}.
\end{thebibliography}

\end{document}
